\chapter{Applications of linear representations}
\label{chap-applications-linear-representations}
\chapteroverview{Character tables turn geometric group actions into concrete
algebraic information. They reveal representation multiplicities and normal
subgroups, while matrix coefficients provide coordinates for data indexed
by a noncommutative group. This chapter develops these applications through
cyclic, dihedral, symmetric, and quaternion groups, and a ranking-data example.}

\section{Representations of familiar groups}
\subsection{Reading a character table}
Let $C_1,\ldots,C_r$ be the conjugacy classes, with $C_1=\{1\}$ and
$k_j=|C_j|$. The character table is the square matrix
$\Phi=(\chi_i(C_j))_{i,j}$, with the trivial character in the first row.
Its first column gives the irreducible degrees $d_i=\chi_i(1)$.
Row orthogonality is weighted by class size:
\[
 \sum_j k_j\chi_i(C_j)\overline{\chi_l(C_j)}=|G|\delta_{il}.
\]
\begin{prop}[Column orthogonality]
\label{prop-orthogonalite-columns}
For conjugacy classes $C,D$,
\[
 \sum_{i=1}^r\chi_i(C)\overline{\chi_i(D)}
   =\begin{cases}|G|/|C|&C=D,\\0&C\ne D.\end{cases}
\]
\end{prop}
\begin{proof}
The indicator of $C$ has coefficient
$c_{f_C}(\rho_i)=|C|\overline{\chi_i(C)}/|G|$.
Apply Plancherel to the indicators of $C,D$, whose inner product is
$|C|\delta_{CD}/|G|$, and simplify. Equivalently, use the unitary matrix
$|G|^{-1/2}\Phi\operatorname{diag}(\sqrt{k_j})$.
\end{proof}
Complex conjugation is essential: character values need not be real.

\subsection{Cyclic groups}
For $G=\langle g_0\rangle$ of order $n$, every irreducible is
one-dimensional. With $\omega_n=e^{2\pi\imath/n}$, all characters are
\[
 \chi_j(g_0^k)=\omega_n^{jk},\qquad 0\leq j,k<n.
\]
Thus the table is the Vandermonde matrix $(\omega_n^{jk})$, and its
normalization by $1/\sqrt n$ is unitary. Choosing $g_0$ identifies the
dual with $\ZZ/n\ZZ$; this recovers the ordinary discrete Fourier theory.

\subsection{Dihedral groups}
For $n\geq3$, the dihedral group $D_n$ has order $2n$ and presentation
\[
 D_n=\langle r,s:r^n=s^2=1,\ srs=r^{-1}\rangle.
\]
Here $r$ is a rotation through $2\pi/n$ and $s$ is a reflection of a
regular $n$-gon. Every element has a unique form $s^\epsilon r^k$,
with $\epsilon\in\{0,1\}$ and $0\leq k<n$.

A one-dimensional representation must send both $r$ and $s$ to signs,
with the additional condition $\rho(r)^n=1$. For even $n$ this gives
four characters, and for odd $n$ only the first two:
\[
 \begin{array}{c|cc}
 &r^k&sr^k\\\hline
 \psi_1&1&1\\\psi_2&1&-1\\
 \psi_3&(-1)^k&(-1)^k\\\psi_4&(-1)^k&(-1)^{k+1}
 \end{array}.
\]
These columns describe two families of elements, not individual conjugacy
classes. The last two rows are used only for even $n$.
For the two-dimensional representations, set
\[
 \rho_h(r)=\begin{pmatrix}\omega_n^h&0\\0&\omega_n^{-h}\end{pmatrix},
 \qquad \rho_h(s)=\begin{pmatrix}0&1\\1&0\end{pmatrix}.
\]
They satisfy the defining relations, and in particular
\[
 \rho_h(sr^k)=\begin{pmatrix}0&\omega_n^{-hk}\\\omega_n^{hk}&0\end{pmatrix}.
\]
Swapping coordinates intertwines $\rho_h$ with $\rho_{n-h}$.
At $h=0$ the representation splits as $\psi_1\oplus\psi_2$; for even
$n$, the case $h=n/2$ splits as $\psi_3\oplus\psi_4$.
For
\[
 1\leq h\leq\lfloor(n-1)/2\rfloor,
\]
the eigenvalues of $\rho_h(r)$ are distinct. Its only invariant lines
are the coordinate axes, which $\rho_h(s)$ exchanges. These representations
are therefore irreducible, with characters
\[
 \chi_h(r^k)=2\cos(2\pi hk/n),\qquad \chi_h(sr^k)=0.
\]
The listed values of $h$ give distinct characters, since the frequencies
$\{h,-h\}$ are distinct modulo $n$. Finally, the degree-squares sums are
$4+4(n/2-1)=2n$ for even $n$ and $2+4(n-1)/2=2n$ for odd $n$.
The lists are complete.

\subsection{The symmetric group on four letters}
\label{sect2-group-s4}
Conjugacy is determined by cycle type. In the order used below, the
five classes have sizes $1,6,8,6,3$: identity, transpositions, three-cycles,
four-cycles, and double transpositions. Let $V_s$ be the standard
three-dimensional representation and $\epsilon$ the sign. Since a
permutation character counts fixed letters, subtracting the trivial
character gives
$\chi_s=(3,1,0,-1,-1)$. Its sign twist, denoted $W$, has character
$(3,-1,0,1,-1)$ and is inequivalent to $V_s$.
Together with the trivial and sign representations, these account for
$1+1+9+9=20$ of the degree-squares sum $24$. There are five irreducibles,
so the remaining one, $W'$, has degree two. The regular-character identity
determines its row:
\[
 \begin{array}{c|rrrrr}
 \text{class size}&1&6&8&6&3\\
 &1&(12)&(123)&(1234)&(12)(34)\\\hline
 \chi_1&1&1&1&1&1\\
 \chi_\epsilon&1&-1&1&-1&1\\
 \chi_s&3&1&0&-1&-1\\
 \chi_W&3&-1&0&1&-1\\
 \chi_{W'}&2&0&-1&0&2
 \end{array}.
\]
Each row has weighted squared norm one, and the rows are mutually orthogonal.
The representation $W'$ is the pullback of the standard representation
of $\mathfrak S_3$ under a quotient map, as the exercises show.

\paragraph{The cube.}
The orientation-preserving symmetries of a cube act on its four body
diagonals. This action identifies the group with $\mathfrak S_4$.
There are $24$ such rotations: an oriented orthonormal coordinate frame
can be sent to any signed coordinate frame of determinant one. The diagonal
action is faithful; a rotation fixing all four diagonal lines can only be
the identity. For example, this follows from the fact that three independent
such directions would be eigenvectors, while their pairwise nonorthogonality
forces the corresponding real eigenvalues to agree; $-I$ reverses orientation.

A half-turn about an axis through opposite edge midpoints has cycle type
$(12)$ on the diagonals. A third-turn about a body diagonal has type
$(123)$; quarter-turns and half-turns about face-center axes have types
$(1234)$ and $(12)(34)$. On the six faces, these fix respectively
$0,0,2,2$ faces, giving
\[
 \chi_{\mathrm{faces}}=(6,0,0,2,2)=\chi_1+\chi_W+\chi_{W'}.
\]
The three-dimensional spatial rotation representation itself is $W$:
its trace is $1+2\cos\theta$, yielding $(3,-1,0,1,-1)$.
The actions on vertices and edges are treated in
Exercise~\oldref{exo-action-face-cube}.

\section{Normal subgroups from character values}
\subsection{The kernel of a character}
\begin{prop}
\label{kernel-prop-characters}
Let $\rho$ be a complex representation of degree $d$. If $g$ has order
$k$, then $\rho(g)$ is diagonalizable with $k$th roots of unity as
eigenvalues. Consequently,
\[
 |\chi_\rho(g)|\leq d,\qquad
 K_{\chi_\rho}:=\{g:\chi_\rho(g)=d\}=\ker\rho.
\]
In particular, $K_{\chi_\rho}$ is a normal subgroup.
\end{prop}
\begin{proof}
The minimal polynomial divides $X^k-1$, which has distinct complex roots.
The trace is therefore a sum of $d$ numbers of modulus one. Its modulus
is at most $d$. If the trace equals the real number $d$, equality in
the triangle inequality forces all eigenvalues to equal one, hence
$\rho(g)=I$. The converse is immediate.
\end{proof}
The weaker equality $|\chi_\rho(g)|=d$ only says that $\rho(g)$ is scalar;
it need not be the identity.

\begin{lem}[Representations of a quotient]
\label{lem-extension-repr-grpe-distignue}
For $N\lhd G$, composition with $\pi:G\to G/N$ identifies representations
of $G/N$ with representations of $G$ that are trivial on $N$. Invariant
subspaces and irreducibility are preserved.
\end{lem}
\begin{proof}
A representation of the quotient pulls back as $\widetilde\rho(g)=\rho(\pi(g))$.
Conversely, if $N\subset\ker\widetilde\rho$, its value depends only on
the coset, defining a representation of $G/N$. Surjectivity of $\pi$
makes the invariant-subspace conditions identical.
\end{proof}

\subsection{Recovering normal subgroups}
\begin{prop}
Every normal subgroup is an intersection of kernels of irreducible
characters, and every such intersection is normal. The empty intersection
is understood to be $G$.
\end{prop}
\begin{proof}
For $N\lhd G$, pull back the regular representation of $G/N$. Since
that representation is faithful, its pullback has kernel exactly $N$.
Decompose it into irreducibles. A direct-sum matrix is the identity
exactly when each of its nonzero summand blocks is the identity, so its
kernel is the intersection of the kernels of the irreducibles that occur.
The converse follows because kernels and their intersections are normal.
\end{proof}
\begin{cor}[A simplicity criterion]
A nontrivial finite group is simple exactly when each nontrivial irreducible
representation is faithful. Equivalently,
\[
 \chi_i(g)\ne\chi_i(1)\quad
 \text{for every nontrivial irreducible $i$ and every $g\ne1$}.
\]
\end{cor}
\begin{proof}
In a simple group, a nontrivial irreducible has a proper normal kernel,
which must be the identity subgroup. Conversely, a proper nontrivial
normal subgroup is an intersection of irreducible kernels and cannot
arise if every nontrivial irreducible kernel is trivial.
\end{proof}
For $\mathfrak S_4$, the nontrivial proper normal subgroups are
\[
 V_4=\{1,(12)(34),(13)(24),(14)(23)\}
       \subset\mathfrak A_4.
\]
They are $\ker\chi_{W'}$ and $\ker\chi_\epsilon$, respectively. The
subgroup generated by the \emph{single} double transposition $(12)(34)$
has only two elements and is not $V_4$; its normal closure is $V_4$.

\section{Spectral analysis on a noncommutative group}
\subsection{Matrix-entry orthogonality}
Choose an invariant Hermitian inner product and an orthonormal basis in
each irreducible space $V_k$, of dimension $d_k$. Write the unitary
representation matrices as $R_k(g)=(r_{ab}^k(g))$. Each entry is a
complex-valued function on $G$.
\begin{thm}[Schur orthogonality for matrix entries]
For valid indices in the respective representations,
\[
 \dotp{r_{ab}^k}{r_{cd}^l}
    =\frac{\delta_{kl}\delta_{ac}\delta_{bd}}{d_k}.
\]
Consequently $(\sqrt{d_k}\,r_{ab}^k)_{k,a,b}$ is an orthonormal basis
of $\CC[G]$ for the normalized inner product.
\end{thm}
\begin{proof}
Average a linear map $T:V_l\to V_k$:
\begin{equation}
\label{eq-orth-functions-coord-interm}
 \widetilde T=\frac1{|G|}\sum_{g\in G}R_k(g)T R_l(g)^*.
\end{equation}
It is zero for $k\ne l$. For $k=l$, Schur's lemma makes it scalar,
and preservation of trace gives $\widetilde T=(\operatorname{tr}T/d_k)I$.
Take $T=E_{bd}$, the matrix with its only nonzero entry in row $b$,
column $d$, and compare entry $(a,c)$. The left side is exactly the
inner product in the assertion. This proves orthogonality and the norm.
There are $\sum_k d_k^2=|G|$ entries, proving completeness.
\end{proof}
Characters are sums of diagonal entries, so their orthogonality is
consistent with this stronger result. The character is independent of
basis; individual entries depend on the chosen orthonormal basis.

\subsection{Fourier coordinates and matrix blocks}
\label{notation-97}
Set $\Delta_{kab}=\sqrt{d_k}\,r_{ab}^k$. For any $f\in\CC[G]$,
\[
 c_f(k,a,b)=\dotp{f}{\Delta_{kab}},\qquad
 f=\sum_{k,a,b}c_f(k,a,b)\Delta_{kab},\qquad
 \norm f_2^2=\sum_{k,a,b}|c_f(k,a,b)|^2.
\]
The matrix-valued transform of the preceding chapter uses $R_k(g)$
without conjugation. Its exact relation to these coordinates is
\begin{equation}
\label{eq-lien-calcul-tfnoncommutative-coeff-fourier}
 (\Ff f(\rho_k))_{ab}
   =\sum_{g\in G}f(g)r_{ab}^k(g)
   =\frac{|G|}{\sqrt{d_k}}\overline{c_{\overline f}(k,a,b)}.
\end{equation}
Both the conjugation and the factor $|G|$ are needed. Equivalently,
conjugated matrix entries provide a Fourier basis directly matching this
transform convention. Plancherel becomes
\[
 \sum_{g\in G}|f(g)|^2
    =\frac1{|G|}\sum_k d_k\norm{\Ff f(\rho_k)}_{\mathrm F}^2,
\]
where $\norm M_{\mathrm F}^2=\sum_{a,b}|M_{ab}|^2$.
A unitary basis change within an irreducible representation conjugates its
Fourier block and preserves this block's energy, although individual
entries and their interpretation change.

A direct evaluation costs $O(|G|^2)$ operations because there are
$|G|$ matrix entries, each defined by a sum of $|G|$ terms. Structured
algorithms exploit subgroup chains and compatible representation bases;
\cite{rockmore-generalized} surveys such methods. Orthogonality guarantees
stable analysis and synthesis, but does not by itself make a signal sparse.
For data analysis, the group structure and basis should reflect meaningful
symmetries of the data.

\subsection{An example with rankings}
Following the statistical viewpoint of \cite{diaconis-representations-statistics},
consider rankings of city, suburb, and countryside. We make the encoding
explicit: $\sigma(j)$ is the \emph{rank assigned to item $j$}, where
items $1,2,3$ are city, suburb, and countryside, and rank one is best.
Thus $(2,3,1)$ means countryside first, city second, suburb third.
The data table retained from the earlier edition is
\[
 \begin{array}{ccc|r}
 \text{city rank}&\text{suburb rank}&\text{countryside rank}&\text{count}\\\hline
 1&2&3&242\\2&1&3&170\\3&2&1&359\\1&3&2&28\\2&3&1&12\\3&1&2&628
 \end{array}.
\]
Its total is $1439$. The most frequent row, $(3,1,2)$, ranks suburb first
and city last. Regard the counts as a function $f:\mathfrak S_3\to\NN$.
Here triples are one-line permutation notation, not cycle notation.

Besides the trivial and sign representations, use the standard representation
on the sum-zero plane in $\CC^3$. In the orthonormal basis
\[
 u=(e_1-e_2)/\sqrt2,\qquad v=(e_1+e_2-2e_3)/\sqrt6,
\]
let $R(\sigma)=B^{\mathsf T}P_\sigma B$, where $B=(u,v)$ and
$P_\sigma e_j=e_{\sigma(j)}$. In table order, the matrices are
\[
\begin{array}{c|c@{\qquad}c|c}
 (1,2,3)&\begin{pmatrix}1&0\\0&1\end{pmatrix}&
 (2,1,3)&\begin{pmatrix}-1&0\\0&1\end{pmatrix}\\[3mm]
 (3,2,1)&\tfrac12\begin{pmatrix}1&-\sqrt3\\-\sqrt3&-1\end{pmatrix}&
 (1,3,2)&\tfrac12\begin{pmatrix}1&\sqrt3\\\sqrt3&-1\end{pmatrix}\\[3mm]
 (2,3,1)&\tfrac12\begin{pmatrix}-1&-\sqrt3\\\sqrt3&-1\end{pmatrix}&
 (3,1,2)&\tfrac12\begin{pmatrix}-1&\sqrt3\\-\sqrt3&-1\end{pmatrix}
\end{array}.
\]
They satisfy the usual permutation composition law. Using $r_{ab}$ for
their entry functions, the exact orthogonal expansion is
\[
 f=\frac16\left(1439\,1+325\epsilon-109r_{11}
          +285\sqrt3\,r_{12}-947\sqrt3\,r_{21}-203r_{22}\right).
\]
For each two-dimensional entry, its squared norm is $1/2$, which accounts
for the factor two relative to an ordinary normalized inner product.
The largest nonconstant coordinate in the normalized basis is the one
associated with
\[
 r_{21}=\frac{\sqrt3}{2}(0,0,-1,1,1,-1).
\]
This function groups rankings according to the item ranked last: it is
negative for city, positive for suburb, and zero for countryside. Those
three totals are $987$, $40$, and $412$, respectively. Its coefficient
therefore measures the strong imbalance between placing city and suburb
last. This interpretation depends on the explicitly chosen rank encoding
and basis; it is not a general property of every matrix coefficient.

\section{Exercises}
\label{sect1-chap6-exercises}

\begin{exo}[Orthogonality of a character table]
\label{exo-orthogonalite-caracteres}
Let $\Phi=(\chi_i(C_j))$ and $K=\operatorname{diag}(|C_1|,\ldots,|C_r|)$.
Show that $\Phi K\Phi^*=|G|I$, and hence that
$U=|G|^{-1/2}\Phi K^{1/2}$ is unitary. Deduce column orthogonality.
The class-size matrix multiplies on the right because classes index columns.
\end{exo}

\begin{exo}[A spatial representation of the dihedral group]
\label{exo-repr-groupe-diedral}
For $n\geq3$, realize $D_n$ by rotations about the $z$-axis through
angles $2\pi k/n$ and half-turns about lines in the $xy$-plane making
angles $k\pi/n$ with the $x$-axis. Show that these transformations form
a representation in $SO_3(\RR)$. Compute its character and decompose
its complexification into irreducibles.
A half-turn about a line is a three-dimensional rotation, whose restriction
to the $xy$-plane is a reflection.
\end{exo}

\begin{exo}[Representations of the symmetric group on three letters]
\label{exo-representation-s3}
Label the triangle vertices $z_1=1$, $z_2=e^{2\pi\imath/3}$, and
$z_3=e^{-2\pi\imath/3}$. In the real basis $(1,\imath)$ of its plane,
compute the matrices of $(12)$ and $(123)$; also identify the reflection
$(23)$. Deduce the full character table of $\mathfrak S_3$.
Here parentheses denote cycles, unlike the one-line triples in the ranking
example.
\end{exo}

\begin{exo}[Actions on cube vertices and edges]
\label{exo-action-face-cube}
The cube rotation group acts on its eight vertices and twelve edges.
Compute both permutation characters in the class order of
Section~\ref{sect2-group-s4}, and use the character table to decompose
the representations. Edges are considered as unordered vertex pairs.
\end{exo}

\begin{exo}[The two-dimensional representation of the symmetric group]
\label{exo-caracteres-s4}
Let $W'$ be the irreducible representation of $\mathfrak S_4$ with
character $(2,0,-1,0,2)$ in the established class order.
\begin{enumerate}
\item Show that every double transposition acts as the identity.
\item Recall why representations trivial on a normal subgroup $H$
are precisely pullbacks of representations of $G/H$.
\item Let $H=\{1,(12)(34),(13)(24),(14)(23)\}$. Use the action on the
three partitions of four letters into two unordered pairs to show
$\mathfrak S_4/H\simeq\mathfrak S_3$.
\item Deduce that $W'$ is the pullback of the standard representation
of $\mathfrak S_3$.
\end{enumerate}
\end{exo}

\begin{exo}[Two-dimensional representations and simple groups]
\label{exo-representation-gpe-simple}
Let $G$ be a finite nonabelian simple group. Prove that it has no
irreducible complex representation of degree two.
\begin{enumerate}
\item Prove Cauchy's theorem: if prime $p$ divides $|G|$, some element
has order $p$. Use cyclic rotation on
\[
 X=\{(g_1,\ldots,g_p)\in G^p:g_1\cdots g_p=1\}.
\]
Check that this subset is rotation-invariant, even for a nonabelian group.
\item If a degree-two irreducible $\rho$ existed, use divisibility of
irreducible degrees to find an involution $t\in G$.
\item Prove that $\rho$ is faithful and takes values in $SL_2(\CC)$.
Show that $\rho(t)=-I$, and derive a contradiction from the center of $G$.
\end{enumerate}
\end{exo}

\begin{exo}[Quaternions and a Fourier basis on eight points]
\label{exo-grpe-quaternionique}
The real quaternion algebra $\mathbb H$ is the four-dimensional real
vector space with basis $(1,i,j,k)$ and bilinear multiplication satisfying
\[
 i^2=j^2=k^2=-1,\qquad ij=k=-ji,\quad jk=i=-kj,\quad ki=j=-ik.
\]
The group $H_8=\{\pm1,\pm i,\pm j,\pm k\}$ consists of eight units.
The group algebra $\RR[H_8]$ has dimension eight and must be distinguished
from $\mathbb H$.
\begin{enumerate}
\item For $q=a+bi+cj+dk$, verify the multiplication table through the
faithful real matrix model
\[
 L(q)=\begin{pmatrix}
 a&-b&-c&-d\\b&a&-d&c\\c&d&a&-b\\d&-c&b&a
 \end{pmatrix}.
\]
Deduce associativity and a real representation of $H_8$. Describe the
surjective algebra map $\RR[H_8]\to\mathbb H$ and its kernel.
\item Verify the real-algebra embedding
\[
 M(q)=\begin{pmatrix}a+\imath b&-c-\imath d\\c-\imath d&a-\imath b\end{pmatrix}.
\]
Compute $M(q)^*M(q)$ and deduce a unitary complex representation of
$H_8$ of dimension two. Prove its irreducibility.
\item Find the four one-dimensional representations and prove that
these five irreducibles form a complete list. Give an explicit orthonormal
basis of $\CC[H_8]$ in the order $(1,-1,i,-i,j,-j,k,-k)$.
\item Transport the group law to $\{0,1\}^3$ by a fixed bijection,
then use the product group $H_8^n$ to construct a Fourier basis on
$\{0,1\}^{3n}$. Specify the dimension $8^n$ and the conjugations and
normalization needed to compute coefficients by tensor products.
\end{enumerate}
This gives an alternative to the abelian Walsh basis. Connections with
learning Boolean functions are discussed in \cite{boneh-learning-boolean}.
\end{exo}

\begin{exo}[Polynomial invariants of cube rotations]
\label{exo-ring-invariants}
Let $G$ be the orientation-preserving symmetry group of the cube with
vertices $(\pm1,\pm1,\pm1)$, acting on $\RR[X,Y,Z]$.
\begin{enumerate}
\item Prove that $X^2+Y^2+Z^2$ is invariant.
\item Show that a rotation is a signed coordinate permutation and sends
$XYZ$ to either $XYZ$ or $-XYZ$. Deduce invariance of $(XYZ)^2$.
\item Define
\[
 \begin{aligned}
 F&=(X+Y+Z)(X+Y-Z)(X-Y+Z)(-X+Y+Z),\\
 Q&=(X^2-Y^2)(X^2-Z^2)(Y^2-Z^2).
 \end{aligned}
\]
Show that $F^2$ and $Q^2$ are invariant. In fact, show that $F$ itself
and the degree-nine polynomial $XYZ\,Q$ are invariant. Interpret the
linear factors through planes associated with the cube's diagonals.
\end{enumerate}
\end{exo}

\begin{exo}[Self-dual binary codes and invariant polynomials]
\label{exo-codes-autoduaux}
Let $\Cc\subset\FF_2^n$ be a self-dual code of dimension $k$. Its
integer weight enumerator is viewed as a polynomial in $\CC[X,Y]$.
\begin{enumerate}
\item Prove $n=2k$.
\item Use homogeneity and the MacWilliams identity to prove invariance
under
\[
 A=\frac1{\sqrt2}\begin{pmatrix}1&1\\1&-1\end{pmatrix}.
\]
\item For $G_1=\langle A\rangle$, show that
\[
 \CC[X,Y]^{G_1}=\CC[X+(\sqrt2-1)Y,\ Y(X-Y)].
\]
Explain why these two generators are algebraically independent.
\item Self-orthogonality makes all codeword weights even. Deduce invariance
under $D=\operatorname{diag}(1,-1)$ and identify
$G_2=\langle A,D\rangle$ as a dihedral group of order sixteen.
The transformation $-I$ alone only encodes even total degree, a weaker
condition than even word weights.
\item Use exact Python Reynolds averages to compute a generating set.
Then show that the two algebraically independent invariants
\[
 P=X^2+Y^2,\qquad Q=X^2Y^2(X^2-Y^2)^2
\]
generate $\CC[X,Y]^{G_2}$. You may use the complex coordinates
$z=X+\imath Y$, $w=X-\imath Y$, or the Molien series after verifying
independence. Deduce the form of a homogeneous self-dual weight enumerator,
and recover the extended Hamming enumerator as $P^4-4Q$.
\end{enumerate}
For larger groups, Reynolds averaging up to Noether's bound can be expensive;
Gr\"obner-basis methods provide additional tools \cite{cox}.
\end{exo}
